\chapter{Problem 3 -- Autentikasi dan Kepercayaan antar-Service}

\begin{ringkas}
Setiap domain kepercayaan memiliki kredensial sendiri yang diterbitkan dan diperiksa oleh pemiliknya, dan tidak ada satu
pun rahasia yang dipakai di dua domain. Kunci BMKG tidak membuka PVMBG, kunci PVMBG tidak membuka BMKG, dan \emph{token} klien
tidak membuka keduanya. Untuk klien, Auth Service menerbitkan JWT berumur 60~detik untuk tiap peran, dan Client-Facing
API menegakkan hak akses \emph{di sisi server}: Media tidak pernah menerima \emph{field} Mentah.
\end{ringkas}

\section{Rancangan solusi dan mekanisme}

Ada lima domain kepercayaan, dan masing-masing dibuktikan terpisah pada bagian bukti.

\begin{table}[!htbp]
\centering\small
\begin{tabularx}{\textwidth}{@{}L{2.5cm}L{4.1cm}L{2.7cm}Y@{}}
\toprule
\textbf{Domain} & \textbf{Kredensial} & \textbf{Diperiksa oleh} & \textbf{Jika bocor}\\
\midrule
BMKG & \code{X-BMKG-Key}, kunci opaque & bmkg-\emph{mock} & hanya BMKG\\
PVMBG & \cmd{Authorization: Bearer pvmbg.v1.<klaim>.<HMAC>}, dengan struktur berbeda & pvmbg-\emph{mock} & hanya PVMBG\\
API internal Aggregator & \code{X-Service-Key} & Aggregator & jalur internal saja; tidak membuka BMKG, PVMBG, atau klien\\
Operator \emph{mock} & \code{X-Admin-Token}, satu per \emph{mock} & masing-masing \emph{mock} & \code{/admin/*} pada \emph{mock} itu saja\\
Klien downstream & \code{client_id} dan \code{client_secret} per peran, lalu JWT & Auth Service (penerbit) dan Client-Facing API (pemeriksa) & hanya peran klien itu\\
\bottomrule
\end{tabularx}
\caption{Lima domain kepercayaan dan kredensialnya masing-masing.}
\end{table}

Aggregator memegang \emph{token} PVMBG, tetapi \textbf{tidak memegang rahasia penanda tangannya}. Akibatnya Aggregator
bisa memanggil PVMBG, tetapi tidak bisa mencetak \emph{token} baru atau mengadministrasi PVMBG.

\subsection{Klien downstream: Media, Tim Lapangan, BNPB Pusat}

\begin{itemize}
 \item \textbf{Tiga identitas yang benar-benar berbeda} (\code{AUTH_MEDIA_*}, \code{AUTH_LAPANGAN_*}, \code{AUTH_PUSAT_*}, masing-masing dengan id dan secret sendiri), bukan satu akun dengan \emph{flag} peran di sisi klien. Peran ditentukan Auth Service dari identitas yang berhasil \emph{login}, lalu ditulis ke klaim \code{role} \emph{di dalam token yang ditandatangani}. Klien tidak dapat memilih perannya sendiri.
 \item \textbf{Access \emph{token}} berupa JWT HS256 dengan klaim \code{sub}, \code{role}, \code{jti}, \code{iat}, dan \code{exp}. TTL-nya \textbf{60~detik} (\code{AUTH_ACCESS_TTL_SEC}, dapat diubah).
 \item \textbf{\emph{Refresh} \emph{token}} berupa 256 bit acak yang \textbf{tidak pernah disimpan dalam bentuk asli}. Yang disimpan hanya hash SHA-256 di SQLite milik Auth Service (volume yang hanya dipasang pada \emph{container} itu). \emph{Token} ini sekali pakai: \cmd{POST /v1/refresh} menandai \emph{token} lama sebagai \code{revoked} dalam transaksi yang sama dengan pembacaannya, lalu menerbitkan pasangan \emph{token} baru. TTL-nya \code{AUTH_REFRESH_TTL_SEC} (24~jam).
 \item \textbf{Penegakan di sisi BNPB pada setiap permintaan} (\code{services/client-api}):
  \begin{enumerate}
   \item Tanpa \cmd{Authorization: Bearer} berarti \textbf{401}. \emph{Token} rusak, kedaluwarsa, salah tanda tangan, atau berperan tidak dikenal berarti \textbf{403}.
   \item \textbf{Media} hanya menerima \emph{field} Ringkasan lewat \emph{whitelist} (\code{policy.FilterItem}). \emph{Field} baru yang kelak muncul di HazardEvent tidak otomatis bocor ke Media.
   \item Permintaan Media atas \emph{field} Mentah (\code{fields=latitude}, \code{include=mentah}, dan sejenisnya) \textbf{ditolak dengan 403}. Aturannya juga berbasis \emph{whitelist} (\code{policy.CheckQuery}): setiap nilai \code{fields} harus berupa \emph{field} Ringkasan, dan \code{include} hanya boleh kosong atau \code{summary}. Variasi penulisan (huruf besar, spasi, \code{fields[]}, parameter berulang, daftar dipisah koma, nama yang tidak dikenal) tidak lolos.
   \item \code{/v1/sources/status} hanya untuk Lapangan dan Pusat, sedangkan Media mendapat 403.
  \end{enumerate}
 \item Kredensial, kunci penandatangan, dan konfigurasi dibaca dari \emph{environment variable}. \emph{Service} berhenti dengan pesan yang jelas jika variabelnya kosong (\code{obs.MustEnv}). Repositori hanya memuat \code{.env.example} yang berisi nilai contoh.
\end{itemize}

\begin{figure}[H]
\centering
\resizebox{0.98\textwidth}{!}{%
\begin{tikzpicture}[font=\sffamily\footnotesize,
  actor/.style={draw=accent, fill=accentlight, rounded corners=3pt, thick, minimum width=2.6cm, minimum height=0.8cm, align=center},
  life/.style={gridgray, dashed},
  msg/.style={-{Stealth[length=2.2mm]}, thick},
  ret/.style={-{Stealth[length=2.2mm]}, thick, dashed},
  lab/.style={font=\sffamily\scriptsize, above, inner sep=2pt, align=center}]
\node[actor] (c) at (0,0) {Klien\\(Lapangan)};
\node[actor] (a) at (5.2,0) {Auth\\Service};
\node[actor] (p) at (10.4,0) {Client-Facing\\API};
\node[actor] (g) at (15.0,0) {Aggregator};
\foreach \n in {c,a,p,g} \draw[life] (\n.south) -- ++(0,-10.4);
\draw[msg] (0,-1.5) -- (5.2,-1.5) node[lab,pos=0.5]{\texttt{POST /v1/token}\\(client\_id, secret)};
\draw[ret] (5.2,-2.7) -- (0,-2.7) node[lab,pos=0.5]{access JWT (60 s, role)\\$+$ refresh token};
\draw[msg] (0,-3.9) -- (10.4,-3.9) node[lab,pos=0.5]{\texttt{GET /v1/hazards}, header Bearer access};
\node[draw=okgreen, fill=okbg, rounded corners=2pt, font=\sffamily\scriptsize, align=center] at (10.4,-4.9) {periksa tanda tangan, \texttt{exp}, peran;\\terapkan kebijakan field};
\draw[msg] (10.4,-5.9) -- (15.0,-5.9) node[lab,pos=0.5]{\texttt{GET /internal/hazards}\\\texttt{X-Service-Key}};
\draw[ret] (15.0,-6.7) -- (10.4,-6.7);
\draw[ret] (10.4,-7.3) -- (0,-7.3) node[lab,pos=0.5]{200: Ringkasan (Media) atau penuh (Lapangan, Pusat)};
\node[draw=warnorange, fill=warnbg, rounded corners=2pt, font=\sffamily\scriptsize, align=center] at (5.2,-8.3) {setelah 60 detik token kedaluwarsa:\\permintaan berikutnya mendapat 401};
\draw[msg] (0,-9.4) -- (5.2,-9.4) node[lab,pos=0.5]{\texttt{POST /v1/refresh}\\(refresh token)};
\draw[ret] (5.2,-10.4) -- (0,-10.4) node[lab,pos=0.5]{pasangan token baru;\\refresh lama tidak berlaku};
\end{tikzpicture}}
\caption{Alur autentikasi klien: login, akses, kedaluwarsa, lalu \emph{refresh} tanpa login ulang.}
\end{figure}


\section{Alasan keputusan dan trade-off}

Tabel berikut membandingkan beberapa cara mengamankan akses klien dengan pilihan kami.

\begin{table}[!htbp]
\centering\small
\begin{tabularx}{\textwidth}{@{}L{4.0cm}Y@{}}
\toprule
\textbf{Alternatif} & \textbf{Alasan tidak dipilih}\\
\midrule
Satu kunci untuk semua & Jika kunci Media bocor, data mentah ikut bocor, dan kunci BMKG bisa menembus PVMBG. Ini justru masalah yang harus diselesaikan.\\
API key statis per klien, tanpa \emph{token} & Tidak ada kedaluwarsa dan \emph{refresh}, padahal Tim Lapangan wajib memakai \emph{token} berumur pendek. Kunci yang bocor akan berlaku selamanya.\\
\emph{Token} opaque dengan introspeksi ke Auth Service pada tiap permintaan & Pencabutan bisa seketika, tetapi setiap permintaan klien menambah satu panggilan jaringan dan Auth Service menjadi titik gagal bagi jalur baca. Hal ini bertentangan dengan tujuan Problem~2.\\
JWT asimetris (RS256 atau EdDSA) & Lebih aman, karena Client-Facing API cukup memegang kunci publik sehingga penyerang yang berhasil membobol API tidak bisa mencetak \emph{token}. Namun pengelolaan kuncinya lebih rumit. \textbf{Dicatat sebagai perbaikan untuk M2.}\\
OAuth2 atau Keycloak lengkap & Berlebihan untuk tiga klien \emph{read-only}, dan sulit didemokan dengan TTL 60~detik.\\
\textbf{JWT HS256 dengan \emph{refresh} \emph{token} opaque (dipilih)} & Verifikasi dilakukan lokal (jalur baca tidak bergantung pada Auth Service), TTL pendek, dan \emph{refresh} \emph{token} dapat dicabut per \emph{token}. Sederhana dan mudah dibuktikan.\\
\bottomrule
\end{tabularx}
\caption{Alternatif mekanisme autentikasi klien.}
\end{table}

Pilihan ini punya konsekuensi. \code{AUTH_JWT_SECRET} dipegang oleh \emph{dua} \emph{service} milik BNPB, yaitu Auth Service dan
Client-Facing API. Rahasia ini independen dari kredensial BMKG, PVMBG, kunci layanan, dan \emph{token} admin, sehingga kebocorannya
tidak merambat ke domain lain. Namun jika Client-Facing API berhasil dibobol, penyerang dapat mencetak \emph{token} klien. Karena
itu JWT asimetris dicatat sebagai langkah berikutnya.

\section{Batasan yang disadari}\label{sec:p3-batasan}

\begin{itemize}
 \item Access \emph{token} lama tidak dicabut seketika saat \emph{refresh}. JWT bersifat \emph{stateless}, sehingga \emph{token} yang belum kedaluwarsa tetap berlaku sampai \code{exp} (paling lama 60~detik). Yang dicabut seketika adalah \emph{refresh token} lama (terbukti: jika dipakai ulang, hasilnya 403). Pencabutan seketika menuntut introspeksi atau daftar \code{jti} yang selalu diperbarui.
 \item Tidak ada deteksi pemakaian ulang \emph{refresh} \emph{token}. Satu keluarga \emph{token} tidak dicabut ketika \emph{refresh} \emph{token} yang sudah terpakai muncul lagi.
 \item \code{/v1/token} tidak membatasi jumlah percobaan (tanpa lockout atau \emph{rate limit} per klien). Perbandingan secret memakai waktu konstan.
 \item \emph{Rate limiter} Client-Facing API bersifat global dan dipasang sebelum autentikasi, sehingga banjir permintaan tanpa \emph{token} ikut menghabiskan jatah. Pembatasan per klien setelah autentikasi menjadi perbaikan untuk M2.
 \item \emph{Token} kedaluwarsa dijawab dengan \textbf{403}, bukan 401, mengikuti kontrak yang dibekukan di \code{docs/contracts.md}.
 \item Tidak ada TLS di dalam jaringan Compose. Secret klien bersifat statis di \emph{environment} tanpa perkakas rotasi. SQLite pada Auth Service hanya cocok untuk satu \emph{instance}, sehingga M2 membutuhkan penyimpanan bersama.
\end{itemize}

\section{Bukti penyelesaian}

Setiap kriteria pada spesifikasi dibuktikan oleh demo yang memeriksa hasilnya sendiri atau oleh tes otomatis.

\begin{table}[!htbp]
\centering\small
\begin{tabularx}{\textwidth}{@{}L{4.4cm}Y@{}}
\toprule
\textbf{Kriteria pada spesifikasi} & \textbf{Bukti}\\
\midrule
Kredensial BMKG ditolak PVMBG dan sebaliknya (401/403) & \cmd{make demo-mock-auth}: kunci BMKG sebagai Bearer $\rightarrow$ 403, \emph{token} PVMBG di \code{X-BMKG-Key} $\rightarrow$ 403, \emph{token} yang dirusak $\rightarrow$ 403, tanpa kredensial $\rightarrow$ 401, dan \emph{token} admin silang $\rightarrow$ 403. Setiap \emph{container} hanya memegang rahasianya sendiri.\\
Media hanya menerima Ringkasan, dan permintaan \emph{field} mentah ditolak server & \cmd{make demo-p3} langkah B (isi respons tanpa \code{latitude}, \code{longitude}, \code{attributes}, dan \code{source_ref_id}) dan langkah C (sepuluh bentuk permintaan Mentah $\rightarrow$ 403, permintaan yang hanya meminta Ringkasan $\rightarrow$ 200). Tes: \code{TestMediaNeverReceivesMentahOnListOrGet} dan \code{TestMediaExplicitMentahRequestIs403}.\\
\emph{Token} Tim Lapangan kedaluwarsa secara alami, \emph{refresh} tanpa \emph{login} ulang, \emph{token} lama ditolak & \cmd{make demo-p3} langkah F: \emph{token} valid $\rightarrow$ tunggu sampai \code{exp} $\rightarrow$ 403 $\rightarrow$ \cmd{POST /v1/refresh} $\rightarrow$ 200 dengan \emph{token} baru \emph{tanpa} \code{client_secret} $\rightarrow$ \emph{refresh} \emph{token} lama dipakai ulang $\rightarrow$ 403.\\
Tidak ada kredensial yang ter-hardcode atau ter-\emph{commit}, dan \code{.env.example} tersedia & \cmd{make secrets}: \code{.env} tidak pernah ter-\emph{commit}, gitleaks bersih (riwayat \emph{commit} dan \emph{working tree}), tidak ada literal kredensial di kode Go non-test, dan \code{.env.example} memuat seluruh variabel.\\
401 dibedakan dari 403 & \cmd{make demo-p3} langkah A dan \code{TestMissingTokenIs401AndBadTokenIs403}: tanpa \emph{token} menghasilkan 401, sedangkan \emph{token} sampah, kedaluwarsa, kunci lain, tanpa \code{exp}, peran tidak dikenal, dan \code{alg=none} semuanya menghasilkan 403.\\
\bottomrule
\end{tabularx}
\caption{Kriteria Problem~3 dan buktinya.}
\end{table}

\subsection*{Kredensial saling menolak (\cmd{make demo-mock-auth})}
\begin{lstlisting}
[1] BMKG mock (X-BMKG-Key)
   PASS own key                         -> HTTP 200
   PASS no credential                   -> HTTP 401
   PASS wrong key                       -> HTTP 403
   PASS PVMBG bearer token              -> HTTP 403
[2] PVMBG mock (Authorization: Bearer <pvmbg-token>)
   PASS own token                       -> HTTP 200
   PASS BMKG key presented as bearer    -> HTTP 403
   PASS tampered PVMBG token            -> HTTP 403
   claims: {"aud":"pvmbg-api","exp":...,"iss":"pvmbg-geologi","sub":"bnpb-aggregator"}
[5] Each container only holds the secrets of its own trust domain
   PASS bmkg-mock has no PVMBG/internal secret; pvmbg-mock has no BMKG/internal secret;
        the aggregator holds tokens but cannot mint or administer
\end{lstlisting}

\subsection*{Media, Mentah, dan token kedaluwarsa (\cmd{make demo-p3})}
\begin{lstlisting}
[1] PASS no header -> 401        PASS empty Bearer -> 401        PASS garbage JWT -> 403
[2] PASS media list omits all Mentah fields on first item
[3] PASS media ?fields=latitude -> 403         PASS media ?include=mentah -> 403
    PASS media ?fields=Latitude -> 403         PASS media ?fields%5B%5D=longitude -> 403
    PASS media ?include=raw -> 403             PASS media asking only for Ringkasan fields -> 200
[4] PASS media /v1/sources/status -> 403       PASS lapangan /v1/sources/status -> 200
[6] expires_in=60s (configured AUTH_ACCESS_TTL_SEC=60)
    PASS fresh access token works
    PASS expired access token -> 403
    PASS refreshed access works without client_id/secret again
    PASS rotated refresh token cannot be reused (403)
\end{lstlisting}

\section{Jawaban}

\begin{jawaban}[Mekanisme yang dipakai]
Mekanisme terdiri dari dua lapis. \textbf{(1) Antarinstansi:} kredensial BMKG (kunci opaque) dan PVMBG (\emph{token} yang
ditandatangani dengan HMAC dan berstruktur lain, dengan rahasia penandatangan hanya ada di \emph{mock} PVMBG) tidak saling
berlaku, dan masing-masing diperiksa oleh pemiliknya. \textbf{(2) Klien downstream:} Auth Service menerbitkan JWT HS256
berumur 60~detik untuk tiap peran, ditambah \emph{refresh} \emph{token} opaque yang sekali pakai. Client-Facing API memeriksa setiap
permintaan dan menegakkan klasifikasi Ringkasan dan Mentah di sisi server.
\end{jawaban}

\begin{jawaban}[Mengapa tidak membuka celah baru]
\begin{itemize}
 \item \textbf{Pemalsuan \emph{token}.} Algoritma dikunci ke HS256 (\code{alg=none} dan algoritma lain ditolak, dan sudah diuji). Klaim \code{exp} wajib ada dan diperiksa (\emph{token} bertanda tangan tanpa \code{exp} ditolak, dan sudah diuji). Peran harus salah satu dari tiga nilai yang dikenal.
 \item \textbf{Eskalasi peran dari sisi klien.} Peran hanya berasal dari klaim yang ditandatangani. \emph{Header} atau parameter buatan klien (misalnya \code{X-Role}) tidak dibaca. Media tidak punya jalur ke \emph{field} Mentah karena respons dibentuk dengan \emph{whitelist}, bukan dengan membuang \emph{field} yang dikenal.
 \item \textbf{Kebocoran lintas domain.} Rahasia tiap domain terpisah dan tidak dipakai ulang. Aggregator tidak memegang rahasia penandatangan PVMBG (diverifikasi oleh \code{demo-mock-auth}).
 \item \textbf{Pencurian dari penyimpanan.} \emph{Refresh} \emph{token} hanya tersimpan sebagai hash, dan databasenya tidak dapat dijangkau \emph{service} lain. Dampak pencurian \emph{refresh} \emph{token} dibatasi oleh sifatnya yang sekali pakai dan oleh TTL.
 \item \textbf{Jendela penyalahgunaan.} \emph{Token} yang dicuri hanya berguna paling lama 60~detik (access \emph{token}) atau satu kali pakai (\emph{refresh} \emph{token}).
 \item \textbf{Batasan yang diakui} ada di bagian~\ref{sec:p3-batasan}: tidak ada pencabutan access \emph{token} seketika, dan rahasia HS256 dipegang dua \emph{service} BNPB.
\end{itemize}
\end{jawaban}
